\documentclass[11pt,a4paper]{article}
\usepackage[margin=25mm]{geometry}
\usepackage{amsmath,amssymb,amsthm,booktabs}
\usepackage[hidelinks,pdfauthor={zslcz and 6.2Astra (AI assistant)},pdftitle={A logarithmic loss in binary prefix coverage for prime reciprocals}]{hyperref}
\newtheorem{theorem}{Theorem}
\newtheorem*{corollary}{Corollary}
\newcommand{\DP}{D_{\mathbb P}}
\newcommand{\DN}{D_{\mathbb N}}
\title{A logarithmic loss in binary prefix coverage\\for prime reciprocals}
\author{zslcz \quad and \quad 6.2Astra (AI assistant)}
\date{Research note, v0.1.0-draft}
\begin{document}
\maketitle
\begin{abstract}
Fix a real base $\beta>2$ and group denominators $n$ into shells
$\beta^{j-1}<n\le\beta^j$. How many digits can we prescribe in the greedy
expansion of $1/n$, starting at its first nonzero digit, if the prescribed
word begins with $1$ and uses only $0$ and $1$? Allowing all integer
denominators realizes every such word of length $j-1$. Restricting the
denominators to primes reduces the depth of complete coverage by at least
$\log_\beta j-O_\beta(1)$ in every sufficiently large shell. We apply
Fourier decay for a nonlinear image of a self-similar measure to place the
denominator interval of one word inside a run of composite integers.
\end{abstract}

\section{Coverage in a denominator shell}
Write $B=\beta^j$, where $j\ge1$ is an integer. For an integer $B/\beta<n\le B$,
set $r_0=1/n$, $d_k=\lfloor\beta r_{k-1}\rfloor$ and
$r_k=\beta r_{k-1}-d_k$. A zero remainder gives a tail of zeros. The first
nonzero digit is $d_j=\lfloor B/n\rfloor$. For $S=\mathbb N$ or $\mathbb P$
(the primes), let $D_S(j)$ be the largest $m\ge1$ such that every word in
$\{1\}\times\{0,1\}^{m-1}$ occurs as $(d_j,\ldots,d_{j+m-1})$ for some
$n\in S$ in the shell. Each word may use a different denominator. Set
$D_S(j)=0$ if the first digit $1$ is missing. The shell is finite and
each denominator supplies at most one word at each depth. A missing word
also rules out all its extensions.

\begin{theorem}\label{main}
For every real $\beta>2$ there exist constants $C_\beta$ and $j_0(\beta)$
such that every integer $j\ge j_0(\beta)$ satisfies
\[
  \DP(j)\le j-\log_\beta j+C_\beta .
\]
\end{theorem}
\begin{corollary}
For the same bases, $\DN(j)-\DP(j)\ge\log_\beta j-O_\beta(1)$.
In particular, the loss of coverage tends to infinity.
\end{corollary}

\paragraph{The interval for a word.}
Given $w=(1,a_1,\dots,a_{m-1})$, write
$A=1+\sum_{r=1}^{m-1}a_r\beta^{-r}$ and $h=\beta^{1-m}$.
Let the remainder after the last digit range over $[0,1)$. Reversing any
step with digit $a\in\{0,1\}$ gives $(a+r)/\beta<1$ for $0\le r<1$.
Thus the full interval for $B/n$ is $[A,A+h)$, contained in $[1,2)$.
The corresponding denominator interval and its length are
\begin{equation}\label{window}
 I_w=\left(\frac{B}{A+h},\frac{B}{A}\right],\qquad
 |I_w|=\frac{Bh}{A(A+h)}.
\end{equation}
For $m=j-1$ and $j\ge2$ we have $Bh=\beta^2$ and $A(A+h)\le4$.
Each interval therefore has length greater than one and contains an
integer. This gives $\DN(j)\ge j-1$ and proves the corollary from
Theorem~\ref{main}.

\section{Proof of Theorem~\ref{main}}
\paragraph{Fourier decay after taking reciprocals.}
Let $\epsilon_1,\epsilon_2,\dots$ be independent fair bits and put
$U=1+\sum_{r\ge1}\epsilon_r\beta^{-r}$. Its law $\mu$ is the homogeneous
self-similar measure for the two maps $x\mapsto 1+(x-1+a)/\beta$,
$a\in\{0,1\}$, with equal weights. Its support lies in $[1,M]$ with
$M=\beta/(\beta-1)<2$. At the first position where two sequences differ,
the discrepancy exceeds the largest possible difference of the tails.
The coding is injective, so $\mu$ has no atoms. On $[1,M]$ take
$F(u)=1/u$. Extend $F''$ by its positive endpoint values outside this
interval and integrate twice, matching $F(1)$ and $F'(1)$. This gives a
$C^2$ extension to $\mathbb R$ with $F''>0$. Theorem~3.1 of Mosquera and
Shmerkin~\cite{ms} applies to $\nu=F_*\mu$, the law of $1/U$. There
are constants $K,\sigma>0$, depending on $\beta$, such that
\[
 |\widehat\nu(\xi)|\le K|\xi|^{-\sigma}\quad(\xi\ne0),\qquad
 \widehat\nu(\xi)=\int e^{-2\pi i\xi x}\,d\nu(x).
\]

\paragraph{Locating a reciprocal modulo a primorial.}
For every integer $1\le q\le B$ and every interval $J$ on the unit circle,
\begin{equation}\label{arc}
 \Pr\big(\{B/(qU)\}\in J\big)=|J|+O_\beta\big((B/q)^{-\sigma}\big),
\end{equation}
where braces denote fractional part. To see the uniformity in $J$,
the $k$th Fourier coefficient of its indicator has magnitude at most
$1/(\pi |k|)$ for $k\ne0$. The corresponding coefficient of the law
modulo one is $\widehat\nu(kB/q)$. The product is bounded in magnitude
by $(K/\pi)(B/q)^{-\sigma}|k|^{-1-\sigma}$, which is summable over
$k\ne0$. To justify the expansion, first integrate the Fej\'er means of
the indicator. They are uniformly bounded and converge off its endpoints.
The distribution modulo one gives those endpoints zero mass, so the
integrals converge to the probability in~\eqref{arc}.

Choose $0<c<\sigma/(2(1+\sigma))$, set $L=\lfloor c\log B\rfloor$,
and let $q=\prod_{p\le L}p$, a product over primes. Here $\log$ denotes
the natural logarithm. The prime number theorem gives $\log q\sim L$, so
eventually $q\le e^{2L}$ and $q=o(B)$. Set $J=[L/(3q),2L/(3q)]$.
For large $j$ it lies in $[0,1)$, and $q^{1+\sigma}/B^\sigma\to0$ by
the choice of $c$. The error in~\eqref{arc} is therefore $o(|J|)$.
The probability is positive, so some binary sequence and some integer $a$
satisfy
\begin{equation}\label{center}
       aq+L/3\le B/U\le aq+2L/3 .
\end{equation}
Since $B/U>B/2$ and $q=o(B)$, we have $a\ge1$ for large $j$.
Every integer $2\le t\le L$ has a prime divisor that divides $q$.
It also divides $aq+t$ and is strictly smaller than $aq+t$.
Thus $aq+2,\ldots,aq+L$ are all composite.

\paragraph{Trapping a whole prefix interval.}
Take $m=1+\lceil\log_\beta(4B/L)\rceil$ and read the prefix of length $m$
of this same $U$, counting its leading $1$. The omitted tail is at most
$h/(\beta-1)<h$, so $B/U\in I_w$. By~\eqref{window}, $|I_w|\le Bh\le L/4$.
Together with~\eqref{center} this places $I_w$ inside
$[aq+L/12,aq+11L/12]$. Once $L\ge24$, every integer in $I_w$ belongs
to the composite run. The word $w$ is missing at depth $m$, and hence
\[
  \DP(j)\le m-1\le j-\log_\beta L+\log_\beta4+1 .
\]
Substituting $L=\lfloor c j\log\beta\rfloor$ proves the theorem.
\hfill$\square$

\section{Finite examples and scope}
For $\beta=\pi$ and $j=2$, the prefix $10$ permits only the denominators
$8$ and $9$. It occurs among integer reciprocals and has no prime witness.
The accompanying code checks every first missing word in the complete
shells up to denominator $10^8$. Rational enclosures certify the integer
endpoints of each interval, and a sieve tests for a prime witness.
\begin{center}
\begin{tabular}{lrrr}
\toprule
Base & Last complete shell $j$ & $\DN(j)$ & $\DP(j)$\\
\midrule
$\pi$ & 16 & 16 & 12\\
$e$   & 18 & 18 & 13\\
\bottomrule
\end{tabular}
\end{center}
These data check the definitions and implementation. The proof gives no
numerical threshold for these samples and leaves the optimal growth rate
of the loss undetermined.

Related work treats membership of individual reciprocals in missing-digit
sets in integer bases~\cite{kom}. Here coverage is a property of a whole
shell and a finite word length that varies with the shell. The argument
uses the prime number theorem and the published Fourier decay estimate
above. Priority for this particular coverage bound remains undetermined.

\paragraph{Contributions and review status.}
zslcz posed the problem and directed the project. The AI assistant
credited as 6.2Astra assisted
with the derivation, literature search, software, and preparation of the
manuscript. The repository includes a Lean verification of the arithmetic
step that excludes primes from a trapped interval. The full asymptotic
argument is given in the text. This draft has not been independently
peer reviewed.
\begin{thebibliography}{9}
\bibitem{ms} C.~A. Mosquera and P.~Shmerkin,
\href{https://www.acadsci.fi/mathematica/Vol43/vol43pp0823-0834.pdf}
{\emph{Self-similar measures: asymptotic bounds for the dimension and Fourier decay of smooth images}},
Ann. Acad. Sci. Fenn. Math. \textbf{43} (2018), 823--834.
\bibitem{kom} S.~D. Kominers,
\emph{A Fixed-Prime Criterion for Reciprocals in Missing-Digit Sets},
\href{https://arxiv.org/abs/2604.11282}{arXiv:2604.11282} (2026).
\end{thebibliography}
\end{document}
